\ELhold

\ELsection{سرفصل‌ها}{سرفصل‌ها}

\begin{ELunit}

\ELfigure{m04-course-outline}{}

\end{ELunit}

\ELhold

\ELsection{جریان‌های الکتریکی دائم}{جریان‌های الکتریکی دائم}

\begin{ELunit}

\begin{itemize}
\tightlist
\item
  در فصل‌های قبل با مسائل الکتریسیته ساکن ناشی از بارهای الکتریکی ساکن سروکار داشتیم.

  \begin{itemize}
  \tightlist
  \item
    اکنون در این فصل به بارهای متحرک الکتریکی که جریان الکتریکی را ایجاد می‌کنند می‌پردازیم.
  \end{itemize}
\item
  جریان‌های الکتریکی ناشی از حرکت بارهای آزاد به سه دسته تقسیم می‌شوند

  \begin{itemize}
  \tightlist
  \item
    \textbf{جریان‌های هدایتی:} در هادی‌ها و نیمه‌هادی‌ها توسط حرکت رانشی الکترون‌ها یا حفره‌ها بوجود می‌آیند.
  \item
    \textbf{جریان‌های الکترولیتی:} ناشی از حرکت یون‌های مثبت و منفی
  \item
    \textbf{جریان‌های انتقالی:} ناشی از حرکت الکترون‌ها یا یون‌ها در خلاء
  \end{itemize}
\item
  در این فصل توجه خود را به جریان‌های هدایتی که قانون اهم بر آن‌ها حاکم است متمرکز می‌کنیم.
\end{itemize}

\ELfigure{m04-map}{}

\begin{itemize}
\tightlist
\item
  مفهوم \textbf{چگالی حجمی جریان} چیست و رابطه \textbf{چگالی جریان انتقالی} چگونه است؟
\item
  رابطه \textbf{چگالی جریان هدایتی} چگونه است؟
\item
  شکل نقطه ای \textbf{قانون اهم} چگونه است؟
\item
  \textbf{معادله پیوستگی} چیست و چگونه از روی آن \textbf{قانون جریان کرشهف} استخراج می شود؟
\item
  رابطه \textbf{قانون ژول} که نشان دهنده میزان اتلاف توان بر اثر عبور جریان از یک جسم هادی است، چگونه است؟
\end{itemize}

\end{ELunit}

\ELhold

\ELsection{مفهوم چگالی جریان و استخراج قانون اهم}{مفهوم چگالی جریان و استخراج قانون اهم}

\begin{ELunit}

\ELfigure{m04-current-element}{}

\begin{itemize}
\tightlist
\item
  حرکت دائمی تعدادی حامل بار، هر یک با بار \ELim{q} گذرنده از جز کوچک سطحیِ \ELim{\Delta s} با سرعت \ELim{\vect{u}} را در نظر بگیرید
\item
  اگر \ELim{N} تعداد حامل‌های بار در واحد حجم باشد، آنگاه مقدار بار گذرنده از سطح \ELim{\Delta s} در مدت زمان \ELim{\Delta t} برابر است با
\end{itemize}

\ELdisp{\Delta Q\ELbr =Nq\vect{u}\cdot\uvec{n}\,\Delta s\,\Delta t\qquad\ELbr (\mathrm{C})}

\begin{itemize}
\tightlist
\item
  چون جریان نرخ زمانی تغییر بار است داریم
\end{itemize}

\ELdisp{\Delta I\ELbr =\frac{\Delta Q}{\Delta t}\ELbr =Nq\vect{u}\cdot\uvec{n}\,\Delta s\ELbr =Nq\vect{u}\cdot\Delta\vect{s}\qquad\ELbr (\mathrm{A})}

\begin{itemize}
\tightlist
\item
  چگالی جریان حجمی با واحد آمپر بر متر مربع را به صورت زیر تعریف می‌کنیم
\end{itemize}

\ELdisp{\vect{J}\ELbr =Nq\vect{u}\qquad\ELbr (\mathrm{A/m^2})}

\end{ELunit}

\begin{ELimportant}{چگالی جریان انتقالی}

\ELdisp{\vect{J}\ELbr =\rho\vect{u}\qquad\ELbr (\mathrm{A/m^2})}

\end{ELimportant}

\begin{ELjoin}

\begin{itemize}
\tightlist
\item
  بنابراین
\end{itemize}

\ELdisp{\Delta I\ELbr =\vect{J}\cdot\Delta\vect{s}}

\begin{itemize}
\tightlist
\item
  به این ترتیب کل جریان گذرنده از سطح دلخواه \ELim{S} برابر است با
\end{itemize}

\begin{ELimportant}{}

\ELdisp{I\ELbr =\int_S\vect{J}\cdot d\vect{s}\qquad\ELbr (\mathrm{A})}

\end{ELimportant}

\end{ELjoin}

\begin{ELunit}

\begin{itemize}
\tightlist
\item
  جریان‌های هدایتی نتیجه حرکت رانشی حامل‌های بار تحت تأثیر یک میدان الکتریکی اعمال شده می‌باشند.
\item
  در مورد اکثر مواد هادی سرعت رانش متوسط، مستقیماً با شدت میدان الکتریکی متناسب است.

  \begin{itemize}
  \tightlist
  \item
    در مورد هادی‌های فلزی داریم
  \end{itemize}
\end{itemize}

\ELdisp{\vect{u}\ELbr =-\mu_e\vect{E}\qquad\ELbr (\mathrm{m/s})}

\begin{itemize}
\tightlist
\item
  \ELim{\mu_e}: ضریب تحرک الکترونی با واحد \ELim{\mathrm{m^2/V\cdot s}}
\end{itemize}

\ELdisp{\vect{J}\ELbr =-\rho_e\mu_e\vect{E}}

\end{ELunit}

\begin{ELjoin}

\begin{itemize}
\tightlist
\item
  \ELim{\rho_e=-Ne}: چگالی حجمی بار الکترون‌های رانشی بوده و یک کمیت منفی است
\end{itemize}

\begin{ELimportant}{شکل نقطه ای قانون اهم}

\ELdisp{\vect{J}\ELbr =\sigma\vect{E}\qquad\ELbr (\mathrm{A/m^2})}

\end{ELimportant}

\end{ELjoin}

\begin{ELunit}

\begin{itemize}
\item
  \ELim{\sigma=-\rho_e\mu_e}: یک پارامتر اساسی محیط بوده و رسانندگی نام دارد
\item
  واحد رسانندگی \ELim{\mathrm{A/V\cdot m}} یا زیمنس بر متر است. عکس رسانندگی را ضریب مقاومت می‌نامند و واحد آن \ELim{\Omega\cdot\mathrm{m}} است.
\item
  بدست آوردن رابطه بین ولتاژ و جریان یک قطعه از ماده‌ای همگن با رسانندگی \ELim{\sigma}، طول \ELim{\ell} و سطح مقطع یکنواخت \ELim{S} با استفاده از شکل نقطه‌ای قانون اهم
\end{itemize}

\ELfigure{m04-resistor}{}

\ELdisp{\left.\begin{aligned}V_{12}=E\ell\quad&\Longrightarrow\quad E=\frac{V_{12}}{\ell}\\I=\int_S\vect{J}\cdot d\vect{s}=JS\quad&\Longrightarrow\quad J=\frac IS\end{aligned}\right\}\quad\ELbr \frac IS\ELbr =\sigma\frac{V_{12}}{\ell}}
\ELdisp{V_{12}\ELbr =\left(\frac{\ell}{\sigma S}\right)I\ELbr =RI}

\end{ELunit}

\begin{ELimportant}{}

\ELdisp{R\ELbr =\frac{\ell}{\sigma S}\qquad\ELbr (\Omega)}
\ELdisp{G\ELbr =\frac1R\ELbr =\sigma\frac S\ell\qquad\ELbr (\mathrm{S})}

\end{ELimportant}

\begin{ELunit}

\begin{itemize}
\tightlist
\item
  \ELim{G}: رسانایی یا عکس مقاومت با واحد مهو یا زیمنس
\end{itemize}

\end{ELunit}

\ELhold

\ELsection{معادله پیوستگی و استخراج قانون جریان کرشهف}{معادله پیوستگی و استخراج قانون جریان کرشهف}

\begin{ELunit}

\ELfigure{m04-map-continuity}{}

\begin{itemize}
\tightlist
\item
  اصل بقای بار یکی از اصول موضوعی اساسی فیزیک است.

  \begin{itemize}
  \tightlist
  \item
    بار الکتریکی نمی‌تواند تولید یا نابود شود.
  \end{itemize}
\item
  حجم دلخواه \ELim{V} را که با سطح \ELim{S} احاطه شده است در نظر بگیرید. بار خالص \ELim{Q} درون این ناحیه وجود دارد. اگر جریان خالص \ELim{I} از سطح به سمت بیرون ناحیه عبور کند، بار درون حجم باید با نرخی برابر با جریان کاهش یابد. برعکس اگر جریان خالصی از سطح به سمت درون ناحیه عبور کند، بار درون حجم باید با نرخی برابر با جریان افزایش یابد.
\item
  جریان خارج شونده از ناحیه، کل شار خروجی بردار چگالی جریان از سطح \ELim{S} است
\end{itemize}

\ELdisp{I\ELbr =\oint_S\vect{J}\cdot d\vect{s}\ELbr =-\frac{dQ}{dt}\ELbr =-\frac{d}{dt}\int_V\rho\,dv}
\ELdisp{\int_V\nabla\cdot\vect{J}\,dv\ELbr =-\int_V\frac{\partial\rho}{\partial t}\,dv}

\end{ELunit}

\begin{ELjoin}

\begin{itemize}
\tightlist
\item
  چون معادله فوق باید مستقل از انتخاب \ELim{V} باشد
\end{itemize}

\begin{ELimportant}{معادله پیوستگی}

\ELdisp{\nabla\cdot\vect{J}\ELbr =-\frac{\partial\rho}{\partial t}\qquad\ELbr (\mathrm{A/m^3})}

\end{ELimportant}

\end{ELjoin}

\begin{ELjoin}

\begin{itemize}
\tightlist
\item
  اگر تغییرات زمانی نداشته باشیم
\end{itemize}

\ELdisp{\nabla\cdot\vect{J}\ELbr =0}

\begin{ELimportant}{قانون جریان کرشهف}

\ELdisp{\oint_S\vect{J}\cdot d\vect{s}\ELbr =0}

\end{ELimportant}

\end{ELjoin}

\begin{ELunit}

\begin{itemize}
\tightlist
\item
  پیش از این گفته بودیم که بارهای قرار داده شده درون یک هادی به سطح آن حرکت کرده به گونه‌ایکه در حالت تعادل درون هادی \ELim{\rho=0} است. اکنون می‌خواهیم این بیان را ثابت کنیم
\end{itemize}

\ELdisp{\left.\begin{aligned}\sigma\nabla\cdot\vect{E}&=-\frac{\partial\rho}{\partial t}\\\nabla\cdot\vect{E}&=\rho/\epsilon\end{aligned}\right\}\quad\ELbr \frac{\partial\rho}{\partial t}+\frac\sigma\epsilon\rho\ELbr =0\quad\ELbr \ELbr \Longrightarrow\quad\ELbr \rho\ELbr =\rho_0e^{-(\sigma/\epsilon)t}\qquad\ELbr (\mathrm{C/m^3})}

\begin{itemize}
\tightlist
\item
  \ELim{\rho_0}: چگالی بار اولیه
\item
  معادله فوق بیان می‌کند که چگالی بار در یک نقطه خاص به صورت نمایی با زمان کاهش می‌یابد.
\item
  چگالی اولیه در زمانی برابر با \ELim{\tau=\epsilon/\sigma} به \ELim{1/e} یا \ELim{36.8\%} مقدار خود کاهش می‌یابد.

  \begin{itemize}
  \tightlist
  \item
    در یک هادی خوب مانند مس \ELim{\tau=1.52\times10^{-19}\,\mathrm{s}} است.
  \item
    برای یک عایق خوب این زمان ممکن است ساعت‌ها یا روزها به طول انجامد.
  \end{itemize}
\end{itemize}

\end{ELunit}

\ELhold

\ELsection{اتلاف توان و قانون ژول}{اتلاف توان و قانون ژول}

\begin{ELunit}

\ELfigure{m04-map-joule}{}

\begin{itemize}
\tightlist
\item
  تحت تأثیر یک میدان الکتریکی، الکترون‌های آزاد در هادی حرکت کرده و با اتم‌های موجود در شبکه کریستالی برخورد می‌کنند. بنابراین انرژی از میدان الکتریکی به اتم‌های تحت ارتعاش منتقل می‌گردد.
\item
  کار انجام شده توسط میدان الکتریکی در حرکت بار \ELim{q} تا فاصله \ELim{\Delta\ell} برابر است با
\end{itemize}

\ELdisp{\Delta W\ELbr =\vect{F}\cdot\Delta\vect{\ell}\ELbr =q\vect{E}\cdot\Delta\vect{\ell}}

\begin{itemize}
\tightlist
\item
  این کار متناظر با توان زیر است
\end{itemize}

\ELdisp{p\ELbr =\lim_{\Delta t\to0}\frac{\Delta w}{\Delta t}\ELbr =q\vect{E}\cdot\vect{u}}

\begin{itemize}
\tightlist
\item
  \ELim{\vect{u}}: سرعت حرکت حامل های بار
\end{itemize}

\ELdisp{dp\ELbr =\rho\,dv\,\vect{E}\cdot\vect{u}}
\ELdisp{dP\ELbr =\vect{E}\cdot\vect{J}\,dv}

\end{ELunit}

\begin{ELjoin}

\begin{itemize}
\tightlist
\item
  به ازای حجم مشخص \ELim{V} کل توان تبدیل شده به حرارت برابر است با
\end{itemize}

\begin{ELimportant}{قانون ژول}

\ELdisp{P\ELbr =\int_V\vect{E}\cdot\vect{J}\,dv\qquad\ELbr (\mathrm{W})}

\end{ELimportant}

\end{ELjoin}

\begin{ELjoin}

\begin{itemize}
\tightlist
\item
  در یک هادی با سطح مقطع ثابت داریم
\end{itemize}

\ELdisp{P\ELbr =\int_LE\,d\ell\int_SJ\,ds\ELbr =VI}

\begin{ELimportant}{}

\ELdisp{P\ELbr =I^2R\qquad\ELbr (\mathrm{W})}

\end{ELimportant}

\end{ELjoin}

\begin{ELunit}

\begin{itemize}
\tightlist
\item
  که همان رابطه آشنای توان اهمی است و حرارت تلف شده در مقاومت \ELim{R} را نشان می‌دهد.
\end{itemize}

\end{ELunit}

\ELhold

\ELsection{بررسی شرایط مرزی چگالی جریان}{بررسی شرایط مرزی چگالی جریان}

\begin{ELunit}

\ELfigure{m04-map-boundary}{}

\begin{itemize}
\tightlist
\item
  هنگامیکه جریان به طور مایل از فصل مشترک بین دو محیط با رسانندگی‌های متفاوت عبور می‌کند، بردار چگالی جریان هم در جهت و هم در اندازه تغییر می‌کند.
\item
  معادلات حاکم بر چگالی جریان دائم عبارتند از
\end{itemize}

\end{ELunit}

\textbf{\lr{Governing Equations for Steady Current Density}}

\Needspace{8\baselineskip}
\begin{ELltr}
\begin{longtable}[c]{@{}
  >{\raggedright\arraybackslash}p{(\columnwidth - 2\tabcolsep) * \real{0.5000}}
  >{\raggedright\arraybackslash}p{(\columnwidth - 2\tabcolsep) * \real{0.5000}}@{}}
\toprule\noalign{}
\ELtablehead \begin{minipage}[b]{\linewidth}\raggedright
\lr{Differential Form}
\end{minipage} & \begin{minipage}[b]{\linewidth}\raggedright
\lr{Integral Form}
\end{minipage} \\\noalign{\vskip 4pt}
\midrule\noalign{}
\endhead
\bottomrule\noalign{}
\endlastfoot
\ELim{\nabla\cdot\vect{J}=0} & \ELim{\oint_S\vect{J}\cdot d\vect{s}=0} \\\noalign{\vskip 4pt}
\ELim{\nabla\times\left(\dfrac{\vect{J}}{\sigma}\right)=0} & \ELim{\oint_C\dfrac1\sigma\vect{J}\cdot\dif\vect{\ell}=0} \\\noalign{\vskip 4pt}
\end{longtable}\end{ELltr}


\begin{ELjoin}

\begin{itemize}
\tightlist
\item
  شرایط مرزی به صورت زیرند
\end{itemize}

\begin{ELimportant}{}

\ELdisp{\nabla\cdot\vect{J}\ELbr =0\quad\ELbr \ELbr \Longrightarrow\quad\ELbr  J_{1n}\ELbr =J_{2n}\qquad\ELbr (\mathrm{A/m^2})}
\ELdisp{\nabla\times(\vect{J}/\sigma)\ELbr =0\quad\ELbr \ELbr \Longrightarrow\quad\ELbr \frac{J_{1t}}{J_{2t}}\ELbr =\frac{\sigma_1}{\sigma_2}}

\end{ELimportant}

\end{ELjoin}

\begin{ELunit}

\begin{itemize}
\tightlist
\item
  هنگامیکه جریان دائمی از مرز دو دی‌الکتریک متفاوت با اتلاف عبور می‌کند
\end{itemize}

\ELdisp{J_{1n}\ELbr =J_{2n}\to\sigma_1E_{1n}\ELbr =\sigma_2E_{2n}}
\ELdisp{D_{1n}-D_{2n}\ELbr =\rho_s\to\epsilon_1E_{1n}-\epsilon_2E_{2n}\ELbr =\rho_s}
\ELdisp{\rho_s\ELbr =\left(\epsilon_1\frac{\sigma_2}{\sigma_1}-\epsilon_2\right)E_{2n}\ELbr =\left(\epsilon_1-\epsilon_2\frac{\sigma_1}{\sigma_2}\right)E_{1n}}

\begin{itemize}
\tightlist
\item
  \ELim{E_{1n}=\uvec{n2}\cdot\vect{E}_1} و \ELim{E_{2n}=\uvec{n2}\cdot\vect{E}_2}
\item
  اگر محیط ۲ نسبت به محیط ۱ هادی به مراتب بهتری باشد
\end{itemize}

\ELdisp{\rho_s\ELbr =\epsilon_1E_{1n}\ELbr =D_{1n}}

\end{ELunit}

\begin{ELjoin}

\begin{itemize}
\tightlist
\item
  که همان رابطه شرط مرزی هادی است.
\end{itemize}

\ELnextnum{4-1}

\begin{ELexample}{}

اختلاف پتانسیل \ELim{\mathcal{V}} به یک خازن صفحه ای موازی با مساحت \ELim{S} اعمال شده است. فضای بین صفحات هادی با دو ماده دی الکتریک تلفاتی به ضخامت های \ELim{d_1} و \ELim{d_2}، گذردهی های \ELim{\epsilon_1} و \ELim{\epsilon_2} و رسانایی های \ELim{\sigma_1} و \ELim{\sigma_2} پر شده است. مطلوبست

\begin{itemize}
\tightlist
\item
  الف. شدت میدان الکتریکی در دو محیط
\item
  ب. چگالی جریان در دو محیط
\item
  ج. چگالی بار سطحی روی صفحات هادی و در مرز دو دی الکتریک
\end{itemize}

\ELfigure{m04-ex1}{}

\begin{ELsolution}

\begin{itemize}
\tightlist
\item
  الف.
\end{itemize}

\ELdisp{J_1\ELbr =J_2\quad\ELbr \ELbr \Longrightarrow\quad\ELbr \begin{aligned}\mathcal{V}&=E_1d_1+E_2d_2\\\sigma_1E_1&=\sigma_2E_2\end{aligned}}
\ELdisp{E_1\ELbr =\frac{\sigma_2\mathcal{V}}{\sigma_2d_1+\sigma_1d_2}\qquad\ELbr  E_2\ELbr =\frac{\sigma_1\mathcal{V}}{\sigma_2d_1+\sigma_1d_2}}

\begin{itemize}
\tightlist
\item
  ب.
\end{itemize}

\ELdisp{J_1\ELbr =J_2\ELbr =\sigma_1E_1\ELbr =\sigma_2E_2\ELbr =\frac{\sigma_1\sigma_2\mathcal{V}}{\sigma_2d_1+\sigma_1d_2}}

\begin{itemize}
\tightlist
\item
  ج.
\end{itemize}

\ELdisp{\uvec{n}\cdot\vect{E}_1\ELbr =\frac{\rho_{s1}}{\epsilon_1}\quad\ELbr \ELbr \Longrightarrow\quad\ELbr \rho_{s1}\ELbr =\epsilon_1E_1\ELbr =\frac{\epsilon_1\sigma_2\mathcal{V}}{\sigma_2d_1+\sigma_1d_2}}
\ELdisp{\uvec{n}\cdot\vect{E}_2\ELbr =\frac{\rho_{s2}}{\epsilon_2}\quad\ELbr \ELbr \Longrightarrow\quad\ELbr \rho_{s2}\ELbr =-\epsilon_2E_2\ELbr =-\frac{\epsilon_2\sigma_1\mathcal{V}}{\sigma_2d_1+\sigma_1d_2}}
\ELdisp{\rho_{si}\ELbr =\left(\epsilon_1-\epsilon_2\frac{\sigma_1}{\sigma_2}\right)E_{1n}\quad\ELbr \ELbr \Longrightarrow\quad\ELbr \begin{aligned}\rho_{si}&=\left(\epsilon_2\frac{\sigma_1}{\sigma_2}-\epsilon_1\right)\frac{\sigma_2\mathcal{V}}{\sigma_2d_1+\sigma_1d_2}\\&=\frac{(\epsilon_2\sigma_1-\epsilon_1\sigma_2)\mathcal{V}}{\sigma_2d_1+\sigma_1d_2}\qquad(\mathrm{C/m^2})\end{aligned}}

\end{ELsolution}

\end{ELexample}

\end{ELjoin}

\ELhold

\ELsection{محاسبه مقاومت}{محاسبه مقاومت}

\begin{ELunit}

\ELfigure{m04-map-resistance}{}

\begin{itemize}
\item
  نحوه محاسبه مقدار \textbf{مقاومت نشتی} در ساختارهای خازنی چگونه است؟
\item
  نحوه محاسبه مقدار مقاومت ساختارهای مختلف چگونه است؟
\item
  ظرفیت بین دو هادی که توسط یک محیط دی‌الکتریک از هم جدا شده‌اند برابر است با
\end{itemize}

\ELdisp{C\ELbr =\frac QV\ELbr =\frac{\oint_S\vect{D}\cdot d\vect{s}}{-\int_L\vect{E}\cdot\dif\vect{\ell}}\ELbr =\frac{\oint_S\epsilon\vect{E}\cdot d\vect{s}}{-\int_L\vect{E}\cdot\dif\vect{\ell}}}

\begin{itemize}
\tightlist
\item
  انتگرال سطحی روی سطح در برگیرنده هادی مثبت و انتگرال خطی از هادی منفی تا هادی مثبت انجام می‌گیرد.
\item
  اگر محیط دی‌الکتریک با اتلاف باشد، جریانی از هادی مثبت به هادی منفی عبور خواهد کرد. مقاومت بین هادی‌ها (مقاومت نشتی) برابر است با
\end{itemize}

\ELdisp{R\ELbr =\frac VI\ELbr =\frac{-\int_L\vect{E}\cdot\dif\vect{\ell}}{\oint_S\vect{J}\cdot d\vect{s}}\ELbr =\frac{-\int_L\vect{E}\cdot\dif\vect{\ell}}{\oint_S\sigma\vect{E}\cdot d\vect{s}}}

\end{ELunit}

\begin{ELjoin}

\begin{itemize}
\tightlist
\item
  درصورتیکه ضریب گذردهی و رسانندگی دارای وابستگی فضایی مشابهی باشند (\ELim{\epsilon(x,y,z)=k\sigma(x,y,z)}) و یا محیط دی‌الکتریک همگن باشد (\ELim{\epsilon} و \ELim{\sigma} وابستگی به مختصات فضایی نداشته باشند)
\end{itemize}

\begin{ELimportant}{}

\ELdisp{RC\ELbr =\frac CG\ELbr =\frac\epsilon\sigma}

\end{ELimportant}

\end{ELjoin}

\ELnextnum{4-2}

\begin{ELexample}{}

مطلوبست محاسبه مقاومت نشتی در واحد طول

\begin{itemize}
\tightlist
\item
  الف. بین هادی های داخلی و خارجی یک کابل هم محور با شعاع \ELim{a} برای هادی داخلی، شعاع \ELim{b} برای هادی خارجی و محیط دی الکتریک با رسانندگی \ELim{\sigma}
\item
  ب. خط انتقال دو سیمه شامل دو سیم با شعاع \ELim{a} و فاصله \ELim{D} درون محیطی با رسانندگی \ELim{\sigma}
\end{itemize}

\begin{ELsolution}

\begin{itemize}
\tightlist
\item
  الف. ظرفیت در واحد طول یک کابل کواکسیال
\end{itemize}

\ELdisp{C_1\ELbr =\frac{2\pi\epsilon}{\ln(b/a)}\qquad\ELbr (\mathrm{F/m})}

\begin{itemize}
\tightlist
\item
  مقاومت نشتی در واحد طول
\end{itemize}

\ELdisp{R_1\ELbr =\frac\epsilon\sigma\left(\frac{1}{C_1}\right)\ELbr =\frac{1}{2\pi\sigma}\ln\left(\frac ba\right)\qquad\ELbr (\Omega\cdot\mathrm{m})}

\begin{itemize}
\tightlist
\item
  ب. ظرفیت در واحد طول یک خط انتقال دوسیمه
\end{itemize}

\ELdisp{C'_1\ELbr =\frac{\pi\epsilon}{\cosh^{-1}\left(\dfrac{D}{2a}\right)}\qquad\ELbr (\mathrm{F/m})}

\begin{itemize}
\tightlist
\item
  مقاومت نشتی در واحد طول
\end{itemize}

\ELdisp{\begin{aligned}R'_1&=\frac\epsilon\sigma\left(\frac{1}{C'_1}\right)=\frac{1}{\pi\sigma}\cosh^{-1}\left(\frac{D}{2a}\right)\\&=\frac{1}{\pi\sigma}\ln\left[\frac{D}{2a}+\sqrt{\left(\frac{D}{2a}\right)^2-1}\right]\qquad(\Omega\cdot\mathrm{m})\end{aligned}}

\end{ELsolution}

\end{ELexample}

\begin{ELunit}

\begin{itemize}
\tightlist
\item
  در وضعیت‌های خاصی مسائل الکتریسیته ساکن و جریان دائم دقیقاً مشابه هم نیستند.
\item
  \textbf{روش اول} محاسبه مقاومت یک قطعه از ماده هادی بین دو سطح یا دو سر به صورت زیر است

  \begin{itemize}
  \tightlist
  \item
    انتخاب دستگاه مختصات مناسب
  \item
    قرار دادن اختلاف پتانسیل \ELim{V_0} بین دو سر هادی
  \item
    یافتن شدت میدان الکتریکی (اگر ماده همگن باشد، حل معادله لاپلاس و بدست آوردن پتانسیل الکتریکی و سپس محاسبه شدت میدان الکتریکی)
  \item
    یافتن جریان کل
    \ELdisp{I\ELbr =\int_S\vect{J}\cdot d\vect{s}\ELbr =\int_S\sigma\vect{E}\cdot d\vect{s}}
  \item
    \ELim{S}: سطح مقطعی که جریان از آن می گذرد
  \item
    مقاومت برابر است با \ELim{V_0/I}
  \end{itemize}
\item
  \textbf{روش دوم} محاسبه مقاومت یک قطعه از ماده هادی بین دو سطح یا دو سر به صورت زیر است (در صورتیکه \ELim{J} به سادگی از روی \ELim{I} قابل تعیین باشد)

  \begin{itemize}
  \tightlist
  \item
    انتخاب دستگاه مختصات مناسب
  \item
    قرار دادن جریان \ELim{I} بین دو سر هادی
  \item
    یافتن \ELim{J} از روی \ELim{I}
  \item
    محاسبه شدت میدان الکتریکی \ELim{\vect{E}=\vect{J}/\sigma} و اختلاف پتانسیل \ELim{V_0}
    \ELdisp{V_0\ELbr =-\int\vect{E}\cdot\dif\vect{\ell}}

    \begin{itemize}
    \tightlist
    \item
      انتگرالگیری از سر با پتانسیل پایین تا سر با پتانسیل بالا انجام می‌گیرد.
    \end{itemize}
  \item
    مقاومت برابر است با \ELim{V_0/I}
  \end{itemize}
\end{itemize}

\end{ELunit}

\ELnextnum{4-3}

\begin{ELexample}{}

یک ماده هادی با ضخامت یکنواخت \ELim{h} و رسانندگی \ELim{\sigma} را به صورت شکل زیر در نظر بگیرید. مقاومت بین دو وجه انتهایی آن را بدست آورید.

\ELfigure{m04-ex3}{}

\begin{ELsolution}

\begin{itemize}
\tightlist
\item
  انتخاب دستگاه مختصات استوانه ای
\end{itemize}

\ELdisp{V\ELbr =0\quad\ELbr \text{\lr{at}}\quad\ELbr \phi\ELbr =0,\qquad\ELbr  V\ELbr =V_0\quad\ELbr \text{\lr{at}}\quad\ELbr \phi\ELbr =\pi/2}
\ELdisp{\frac{d^2V}{d\phi^2}\ELbr =0\quad\ELbr \ELbr \Longrightarrow\quad\ELbr  V\ELbr =c_1\phi+c_2\quad\ELbr \ELbr \Longrightarrow\quad\ELbr  V\ELbr =\frac{2V_0}{\pi}\phi}
\ELdisp{\begin{aligned}\vect{J}=\sigma\vect{E}&=-\sigma\nabla V\\&=-\uvec{\phi}\sigma\frac{\partial V}{r\,\partial\phi}=-\uvec{\phi}\frac{2\sigma V_0}{\pi r}\end{aligned}}
\ELdisp{\begin{aligned}I=\int_S\vect{J}\cdot d\vect{s}&=\int_0^h\!\!\int_a^b\left(-\uvec{\phi}\frac{2\sigma V_0}{\pi r}\right)\cdot\left(-\uvec{\phi}\,dr\,dz\right)\\&=\frac{2\sigma hV_0}{\pi}\ln\frac ba\end{aligned}}
\ELdisp{R\ELbr =\frac{V_0}{I}\ELbr =\frac{\pi}{2\sigma h\ln(b/a)}}

\end{ELsolution}

\end{ELexample}
